\section*{Задачи}
\addcontentsline{toc}{section}{Задачи}

\begin{enumerate}
    \item \textit{(Демо 2026, №1)} Составить систему линейных уравнений, задающую линейную оболочку векторов 
    \[
    a_1 = (1, 1, 2, 1, 2)^T, a_2 = (0, -1, -2, 1, -1)^T, a_3 = (3, 1, 2, 5, 4)^T
    \]    
    в $\mathbb{R}^5$

    \item \textit{(Проскуряков, №1312)} Найти систему линейных уравнений, задающую линейное подпространство, натянутое на систему векторов:
    \[
    a_1 = (1, -1, 1, 0), \quad a_2 = (1, 1, 0, 1), \quad a_3 = (2, 0, 1, 1)
    \]

    \item \textit{(Проскуряков, №1313)} Найти систему линейных уравнений, задающую линейное подпространство, натянутое на систему векторов:
    \[
    a_1 = (1, -1, 1, -1, 1), \quad a_2 = (1, 1, 0, 0, 3), \quad a_3 = (3, 1, 1, -1, 7), \quad a_4 = (0, 2, -1, 1, 2)
    \]

    \item \textit{(Кострикин, №35.11-а)} Найти базис и размерность линейной оболочки системы векторов в $\mathbb{R}^4$:
    \[
    a_1 = (1, 0, 0, -1), \quad a_2 = (2, 1, 1, 0), \quad a_3 = (1, 1, 1, 1), \quad a_4 = (1, 2, 3, 4), \quad a_5 = (0, 1, 2, 3)
    \]

    \item \textit{(Кострикин, №35.11-б)} Найти базис и размерность линейной оболочки системы векторов в $\mathbb{R}^5$:
    \[
    a_1 = (1, 1, 1, 1, 0), \quad a_2 = (1, 1, -1, -1, -1), \quad a_3 = (2, 2, 0, 0, -1),
    \]
    \[
    a_4 = (1, 1, 5, 5, 2), \quad a_5 = (1, -1, -1, 0, 0)
    \]


    \item \textit{(Демо 2026, №2)} Найти размерности и базисы суммы и пересечения подпространств $V_1, V_2$ в $\mathbb{R}^4$,
    \[
    V_1 = \langle a_1, a_2, a_3 \rangle,\qquad
    V_2 = \langle b_1, b_2 \rangle,
    \]
    где
    \[
    a_1 = (1, 0, -3, -2)^T,\quad
    a_2 = (7, 1, 9, 14)^T,\quad
    a_3 = (-4, 1, 2, -9)^T,
    \]
    \[
    b_1 = (10, 1, 0, 8)^T,\quad
    b_2 = (-3, 0, 1, -3)^T.
    \]

    \item \textit{(Проскуряков, №1317)} Найти размерность $s$ суммы и размерность $d$ пересечения линейных подпространств $L_1 = \langle a_1, a_2 \rangle$ и $L_2 = \langle b_1, b_2 \rangle$ в $\mathbb{R}^4$, где
    \[
    a_1 = (1, 2, 0, 1)^T, \quad a_2 = (1, 1, 1, 0)^T,
    \]
    \[
    b_1 = (1, 0, 1, 0)^T, \quad b_2 = (1, 3, 0, 1)^T.
    \]

    \item \textit{(Проскуряков, №1318)} Найти размерность $s$ суммы и размерность $d$ пересечения линейных подпространств $L_1 = \langle a_1, a_2, a_3 \rangle$ и $L_2 = \langle b_1, b_2, b_3 \rangle$ в $\mathbb{R}^4$, где
        \[
        a_1 = (1, 1, 1, 1)^T, \quad a_2 = (1, -1, 1, -1)^T, \quad a_3 = (1, 3, 1, 3)^T,
        \]
        \[
        b_1 = (1, 2, 0, 2)^T, \quad b_2 = (1, 2, 1, 2)^T, \quad b_3 = (3, 1, 3, 1)^T.
        \]

    \item \textit{(Демо 2026, №3)} Доказать, что пространство $\mathbb{R}^4$ является прямой суммой подпространств
    \[
    L_1=\langle a_1,a_2\rangle,\qquad L_2=\langle b_1,b_2\rangle,
    \]
    и разложить вектор
    \[
    x=(0,-2,2,0)^T
    \]
    на сумму проекций на эти подпространства, где
    \[
    a_1=(1,1,1,0)^T,\quad a_2=(1,1,0,1)^T,\quad
    b_1=(1,0,1,1)^T,\quad b_2=(1,1,-1,-1)^T.
    \]

    \item \textit{(Проскуряков, №1328)} Доказать, что пространство $\mathbb{R}_n$ есть прямая сумма двух линейных подпространств: $L_1$, заданного уравнением $x_1 + x_2 + \dots + x_n = 0$, и $L_2$, заданного системой уравнений $x_1 = x_2 = \dots = x_n$. Найти проекции единичных векторов $e_1 = (1, 0, 0, \dots, 0), e_2 = (0, 1, 0, \dots, 0), \dots, e_n = (0, 0, 0, \dots, 1)$ на $L_1$ параллельно $L_2$ и на $L_2$ параллельно $L_1$.

    \item \textit{(Проскуряков, №1329)} Доказать, что пространство всех квадратных матриц порядка $n$ есть прямая сумма линейных подпространств $L_1$ — симметрических и $L_2$ — кососимметрических матриц. Найти проекции $A_1$ и $A_2$ матрицы
    \[
    A = \begin{pmatrix}
    1 & 1 & \dots & 1 \\
    0 & 1 & \dots & 1 \\
    \vdots & \vdots & \ddots & \vdots \\
    0 & 0 & \dots & 1
    \end{pmatrix}
    \]
    на $L_1$ параллельно $L_2$ и на $L_2$ параллельно $L_1$.

    \item \textit{(Кострикин, №35.19)} Пусть в пространстве $\mathbb{R}^4$ заданы подпространства:
    \[
    U = \langle (1, 1, 1, 1), (-1, -2, 0, 1) \rangle, \qquad V = \langle (-1, -1, 1, -1), (2, 2, 0, 1) \rangle
    \]
    Доказать, что $\mathbb{R}^4 = U \oplus V$, и найти проекцию вектора $x = (4, 2, 4, 4)$ на подпространство $U$ параллельно $V$.

    \item \textit{(Кострикин, №35.15-а)} Найти базисы суммы и пересечения линейных оболочек $\langle a_1, a_2, a_3 \rangle$ и $\langle b_1, b_2, b_3 \rangle$:
    \[
    a_1 = (1, 2, 1), \quad a_2 = (1, 1, -1), \quad a_3 = (1, 3, 3);
    \]
    \[
    b_1 = (1, 2, 2), \quad b_2 = (2, 3, -1), \quad b_3 = (1, 1, -3).
    \]
\end{enumerate}
